\chapter{Solutions to the exercises of Chapters 12 to 16}\markboth{Solutions 12 to 16}{Solutions 12 to 16}
\begin{description}[leftmargin=0pt,itemsep=4pt]
\item[12.1] $d\theta/d\psi$ with $x=\alpha|\psi|$, $dx/d\psi=-\alpha$ for $\psi<0$: $(\theta_s-\theta_r)(-m)(1+x^n)^{-m-1}n\,x^{n-1}(-\alpha)$, which is Equation (12.2).
\item[12.2] $C(-300)=0.352(0.036)(0.359)(1.56)(3.79)(0.00625)=1.7\times10^{-4}$ cm$^{-1}$.
\item[12.3] If $\psi^{k+1}=\psi^k$ the capacity term vanishes and the remaining term $(\theta^k-\theta^j)/\Delta t$ is the storage change with $\theta^k=\theta(\psi^{j+1})$; the right-hand side is the flux divergence at $\psi^{j+1}$, which is Equation (12.4).
\item[12.4] At a sharp front $K$ changes by orders of magnitude between neighbouring cells. The arithmetic mean is dominated by the wet value and lets water into the dry cell far faster than it can go; the geometric mean is close to the smaller value.
\item[12.5] No. A change of 6 per cent shows the result still depends on the grid. Refine again (and reduce $\Delta t$) until successive changes are below about 1 to 2 per cent.
\item[12.6] $(13.81-7.71)/0.25=24.4$ cm\,d$^{-1}$, which is $98$ per cent of $\Ks=24.96$.
\item[13.1] $e_s=4.243$, $e_a=1.485$, $\Delta=0.2434$, $\gamma=0.0674$; $ET_0=9.07$ mm\,d$^{-1}$.
\item[13.2] $1/2.45=0.408$ mm per MJ\,m$^{-2}$\,d$^{-1}$.
\item[13.3] Stage 1 lasts $6/5=1.2$ d; $t-t_1=4.8$ d; $E=6+2.5\sqrt{4.8}=11.5$ mm.
\item[13.4] Each wetting restarts stage 1, where evaporation proceeds at the potential rate; a single large wetting spends most of its time in stage 2, with the dry surface layer slowing evaporation.
\item[14.1] $RAW=0.55(95)=52.3$ mm; $K_s=(95-70)/(0.45\cdot95)=0.585$.
\item[14.2] $1-0.9(0.3)=0.73$.
\item[14.3] Sand: $RAW=30$ mm, $30/5=6$ d. Loam: $RAW=75$ mm, $15$ d.
\item[14.4] A sensor trigger beyond the reach is never read, so the crop is not irrigated even as the stress coefficient falls; the field is rainfed and the sand's small $TAW$ means the stress comes within days.
\item[15.1] $\tau=(60/3)^{1/0.6}=20^{1.667}=147$ min.
\item[15.2] $r=2.0/225=8.89$ mm\,h$^{-1}$; $80/8.89=9.0$ h.
\item[15.3] Mean $9$; deviations $3,1,0,1,3$ sum to $8$; $CU=100(1-8/45)=82.2$ per cent.
\item[15.4] $r=4/(0.75\cdot2)=2.67$ mm\,h$^{-1}$; $5/2.67=1.88$ h.
\item[15.5] $L^2=(8\cdot1\cdot2\cdot0.5+4\cdot1\cdot0.25)/0.007=9/0.007=1{,}286$, $L=35.9$ m.
\item[15.6] Water arrives at the top before the bottom, so the top gets a longer opportunity time than the bottom, and some water runs off the tail.
\item[16.1] Rank differences $-1,1,0,-1,1$, $\sum d^2=4$, $\rho=1-24/120=0.8$.
\item[16.2] Cells within a lattice share a soil and a design, so they are the units that are exchangeable; resampling across different lattices would mix designs.
\item[16.3] $\mathcal L(3.5\times10^{-10})=0$; the predicted reaches are $199,\ 1{,}615$ and $127{,}830$ cm, all inside the intervals. It is on the plateau.
\item[16.4] Because the left-out soil cannot influence the fit, so the error measures prediction and not description.
\item[16.5] $1/10!=1/3{,}628{,}800=2.8\times10^{-7}$.
\end{description}
